% presentation/slides/04_cas1.tex
\begin{frame}{Le Cas d'Application Linéaire 1D}
    \begin{columns}[c]
        \begin{column}{0.55\textwidth}
            \textbf{Modèle Jouet : Décalages circulaires (Translations sur le Tore)}
            \begin{itemize}
                \item $\theta$ est un signal périodique modélisé dans la base de Fourier.
                \item Les variables cachées $Z_i$ représentent des translations uniformes sur la grille discretisée.
                \item L'opérateur $A_{Z_i}$ applique un déphasage en fréquences.
            \end{itemize}
            \vspace{0.2cm}
            \alert{\textbf{Piège Statistique :}} La moyenne empirique brute $\frac{1}{n}\sum Y_i$ converge vers un signal plat. L'information de phase s'efface !
        \end{column}
        \begin{column}{0.45\textwidth}
            \centering
            % Remplacer par \includegraphics[width=0.9\textwidth]{figures/signaux_1D.pdf}
            \framebox{\begin{minipage}{0.9\textwidth}\centering\vspace{1.2cm}\textcolor{DarkBeige}{\textbf{[ figures/signaux_1D.pdf ]}}\\\small Signaux simulés translatés\vspace{1.2cm}\end{minipage}}
        \end{column}
    \end{columns}
\end{frame}